\section{Resultados}

\subsection{Seleção e teste reservado}

Na validação, com 1.205 consultas, BM25 apresentou Top-1 de 60,50\%, e a inclusão dos trigramas elevou esse resultado para 95,19\%. As quatro configurações de Random Forest tiveram Top-1 de 99,42\% e MRR por nome de 0,9942. O desempate selecionou a primeira configuração da grade. Em relação aos trigramas, o modelo acrescentou 52 acertos de primeira posição e perdeu um; portanto, o desempenho na validação não foi perfeito.

No teste reservado, os trigramas elevaram a cobertura média dos relevantes nos 100 candidatos de 86,50\% para 100,00\%. A Tabela~\ref{tab:principal} apresenta os resultados de recuperação e reordenação. BM25 acertou 767 das 1.207 consultas na primeira posição, os trigramas acertaram 1.171 e o Random Forest acertou 1.207. O modelo corrigiu as 36 falhas restantes dos trigramas, sem perdas, aumentando Top-1 em 2,98 pontos percentuais. Não houve alvo ausente dos candidatos dos trigramas nesse teste.

% Gerado por gerar_tabelas.py a partir dos resultados congelados.
\begin{table}[tb]
\caption{Teste reservado: médias por consulta (1.207 consultas).}
\label{tab:principal}
\centering
\small
\setlength{\tabcolsep}{3pt}
\begin{tabular}{lrrr}
\toprule
Método & Top-1 (\%) & R@100 (\%) & MRR@100 \\
\midrule
BM25 & 63,55 & 86,50 & 0,6889 \\
BM25 + trigramas & 97,02 & 100,00 & 0,9808 \\
Damerau--Levenshtein & 91,80 & 100,00 & 0,9264 \\
Jaro--Winkler & 95,36 & 100,00 & 0,9670 \\
Token Sort & 98,18 & 100,00 & 0,9854 \\
Token Set & 98,59 & 100,00 & 0,9923 \\
Random Forest & 100,00 & 100,00 & 1,0000 \\
\bottomrule
\end{tabular}
\par\smallskip{\footnotesize As quatro similaridades isoladas foram comparadas após o teste original, sem modificar o modelo.}
\end{table}


Na comparação adicional das similaridades isoladas, Token Set obteve o maior Top-1, com 1.190 acertos. O Random Forest acrescentou 17 acertos em relação a esse método, diferença de 1,41 ponto percentual. Token Sort, Jaro--Winkler e Damerau--Levenshtein alcançaram, respectivamente, 1.185, 1.151 e 1.108 acertos. Todos os reordenadores conservaram Recall@100 de 100,00\%, pois compartilharam os candidatos dos trigramas.

A média por nome preservou a tendência: Top-1 foi 61,59\% para BM25, 96,89\% para trigramas e 100,00\% para Random Forest. Isso reduz a possibilidade de que a comparação principal seja explicada apenas pela quantidade desigual de variantes por nome, embora não substitua uma avaliação em novos dados.

\subsection{Posições iniciais e tipos de transformação}

A Tabela~\ref{tab:posicoes} distingue a presença de um relevante nas posições iniciais da proporção de relevantes nessas posições. Com Random Forest, Acerto@3 e Acerto@5 foram 100,00\%, enquanto Precision@3 e Precision@5 foram 33,66\% e 20,20\%. Esses valores são compatíveis com consultas que geralmente têm apenas um ID relevante. Não representam uma queda do acerto da primeira sugestão para 20\%.

% Gerado por gerar_tabelas.py a partir dos resultados congelados.
\begin{table}[tb]
\caption{Posições iniciais no teste: acerto (H) e precisão (P), em porcentagem.}
\label{tab:posicoes}
\centering
\small
\setlength{\tabcolsep}{3pt}
\begin{tabular}{lrrrr}
\toprule
Método & H@3 & H@5 & P@3 & P@5 \\
\midrule
BM25 & 72,91 & 76,14 & 24,41 & 15,29 \\
BM25 + trigramas & 99,09 & 99,50 & 33,36 & 20,10 \\
Damerau--Levenshtein & 92,71 & 93,70 & 31,23 & 18,94 \\
Jaro--Winkler & 97,85 & 98,26 & 32,95 & 19,85 \\
Token Sort & 98,59 & 98,92 & 33,20 & 19,98 \\
Token Set & 99,83 & 99,83 & 33,61 & 20,17 \\
Random Forest & 100,00 & 100,00 & 33,66 & 20,20 \\
\bottomrule
\end{tabular}

\end{table}


A Tabela~\ref{tab:tipos} detalha o comportamento por transformação. No conjunto de 167 consultas originais, todos os métodos acertaram a primeira posição. As diferenças apareceram nas variantes. Os trigramas falharam em 17 transposições, nove casos de duas omissões, cinco substituições, quatro omissões e uma abreviação. O Random Forest corrigiu esses casos.

% Gerado por gerar_tabelas.py a partir dos resultados congelados.
\begin{table}[tb]
\caption{Top-1 por transformação no teste (\%).}
\label{tab:tipos}
\centering
\small
\setlength{\tabcolsep}{3pt}
\begin{tabular}{lrrrrr}
\toprule
Transformação & n & BM25 & Tri. & DL & RF \\
\midrule
Original & 167 & 100,00 & 100,00 & 100,00 & 100,00 \\
Normalização & 2 & 100,00 & 100,00 & 100,00 & 100,00 \\
Abreviação & 166 & 98,19 & 99,40 & 98,80 & 100,00 \\
Omissão & 165 & 47,88 & 97,58 & 100,00 & 100,00 \\
Substituição & 165 & 47,27 & 96,97 & 100,00 & 100,00 \\
Inserção & 165 & 46,67 & 100,00 & 100,00 & 100,00 \\
Transposição & 165 & 46,67 & 89,70 & 100,00 & 100,00 \\
Duas omissões & 97 & 9,28 & 90,72 & 100,00 & 100,00 \\
Ordem & 115 & 100,00 & 100,00 & 15,65 & 100,00 \\
\bottomrule
\end{tabular}
\par\smallskip{\footnotesize Tri.: BM25 com trigramas; DL: Damerau--Levenshtein; RF: Random Forest.}
\end{table}


Das 99 falhas de Damerau--Levenshtein, 97 ocorreram na inversão da ordem de palavras e duas em abreviações. Jaro--Winkler apresentou 44 falhas de ordem, nove de abreviação, uma de omissão e duas de duas omissões. Token Set distribuiu suas 17 falhas entre omissão, transposição, substituição, inserção e duas omissões, sem falhar na categoria de ordem. Esses padrões mostram diferenças de sensibilidade entre as representações.

Um exemplo foi a consulta \emph{rua ameli sa barbosa}. Token Set colocou \emph{rua sa barbosa} antes da referência \emph{rua amelia sa barbosa}; o Random Forest colocou a referência em primeiro. Outro caso, \emph{rua cotsa}, teve a referência \emph{rua costa} recuperada na posição 76 pelos trigramas e promovida para a primeira posição pelo modelo. Os exemplos ilustram ganhos de ordenação dentro de candidatos já recuperados.

\subsection{Complemento exploratório}

A Tabela~\ref{tab:complementar} separa as três categorias complementares. Random Forest e Damerau--Levenshtein acertaram as 179 consultas na primeira posição. Nas dez grafias documentadas, os trigramas e todos os reordenadores já atingiram 100,00\%; BM25 atingiu 20,00\%. Esse grupo não demonstrou ganho adicional do modelo sobre os trigramas.

% Gerado por gerar_tabelas.py a partir dos resultados congelados.
\begin{table}[tb]
\caption{Top-1 no complemento exploratório (\%).}
\label{tab:complementar}
\centering
\small
\setlength{\tabcolsep}{3pt}
\begin{tabular}{lrrr}
\toprule
Método & Doc. (10) & Hist. (90) & Fon. (79) \\
\midrule
BM25 & 20,00 & 44,44 & 69,62 \\
BM25 + trigramas & 100,00 & 93,33 & 92,41 \\
Damerau--Levenshtein & 100,00 & 100,00 & 100,00 \\
Jaro--Winkler & 100,00 & 100,00 & 97,47 \\
Token Sort & 100,00 & 98,89 & 89,87 \\
Token Set & 100,00 & 94,44 & 98,73 \\
Random Forest & 100,00 & 100,00 & 100,00 \\
\bottomrule
\end{tabular}
\par\smallskip{\footnotesize Doc.: grafias documentadas; Hist. e Fon.: variantes históricas e fonéticas simuladas.}
\end{table}


Nas variantes históricas simuladas, os trigramas tiveram Top-1 de 93,33\% e MRR@100 de 0,9571. Nas fonéticas simuladas, os valores foram 92,41\% e 0,9513. O Random Forest corrigiu seis falhas de cada categoria, sem perdas. O Recall@100 do BM25 foi 90,00\% no grupo histórico simulado e 100,00\% nos outros dois; trigramas e reordenadores tiveram cobertura completa nas três categorias.

No recorte de 32 consultas novas associadas a nomes do teste original, os trigramas acertaram 30 e o Random Forest 32. Esse resultado permanece exploratório, pois o conjunto foi elaborado depois da avaliação principal. O empate de Damerau--Levenshtein com o modelo no complemento contrasta com sua dificuldade nas inversões de ordem do teste reservado.

\subsection{Correspondência funcional, espacial e latência}

Os rankings completos produzidos pela função utilizada no serviço coincidiram com o processamento offline nas 1.207 consultas do teste, usando o cache de candidatos. Na validação HTTP ao vivo, as cinco sugestões e o identificador do modelo coincidiram com as referências nas 336 buscas medidas. As 112 chamadas de aquecimento também passaram e foram contabilizadas separadamente.

Os 219 controles espaciais concordaram com o backup, incluindo respostas vazias, múltiplos pontos e números decimais. Os 60 fluxos de duas chamadas também passaram. A Tabela~\ref{tab:latencia} apresenta os tempos medidos. A busca teve mediana de 14,97 ms; o fluxo de busca e obtenção dos pontos teve mediana de 20,58 ms e P95 de 46,99 ms.

% Gerado por gerar_tabelas.py a partir dos resultados congelados.
\begin{table}[tb]
\caption{Validação local: verificações e tempos HTTP em ms.}
\label{tab:latencia}
\centering
\small
\setlength{\tabcolsep}{3pt}
\begin{tabular}{lrrrr}
\toprule
Etapa & Corretas & Mediana & P95 & Média \\
\midrule
Busca & 336/336 & 14,97 & 33,09 & 18,85 \\
Pontos & 219/219 & 3,92 & 22,54 & 6,60 \\
Busca + pontos & 60/60 & 20,58 & 46,99 & 27,06 \\
\bottomrule
\end{tabular}
\par\smallskip{\footnotesize Busca: 112 consultas em três rodadas. Fluxos: 20 casos em três rodadas. Aquecimento excluído.}
\end{table}


Não houve falhas ou interrupções nessa coleta. Esses resultados demonstram consistência do protótipo local com suas referências e atendimento rápido nas condições sequenciais avaliadas. Não incluem desenho do mapa nem estimam comportamento sob carga ou em acesso remoto.
